\documentclass[aspectratio=169]{beamer}
\usetheme{simple}

\input{preamble/preamble.tex}
\input{preamble/preamble_math.tex}
% \input{preamble/preamble_acronyms.tex}

\title{Mixture Models and EM}
\subtitle{STATS 305B: Applied Statistics II}
\author{Scott Linderman}
\date{\today}



\begin{document}

\maketitle


\begin{frame}{Recap}

We have a lot of tools in our toolkit now! We learned about exponential family distributions, which form the building blocks of more complex models. 
    
We also learned about Bayesian inference algorithms, like MCMC and VI, to infer posterior distributions of more complex models

Today we'll start learning about \textbf{latent variable models} for complex datasets. We'll start with the simplest latent variable model --- \textbf{mixture models}.

\end{frame}


\begin{frame}{Outline}

    
\begin{itemize}
    \item Specifying a mixture model
    \item K-Means and MAP estimation
    \item Expectation Maximization (EM) and MLE
    \item Connecting EM and VI
\end{itemize}

\end{frame}

    



\begin{frame}[allowframebreaks]{Motivation}\label{motivation}
    \begin{figure}
        \centering
        \includegraphics[width=0.65\textwidth]{../lectures/figures/scrnaseq.jpg}
        \caption{Single Cell RNA Sequencing. \emph{From \{cite:t\}\texttt{Kiselev2019-bt}.}}
    \end{figure}
\end{frame}

\begin{frame}[allowframebreaks]{Motivation}\label{motivation}
\begin{figure}
\centering
\includegraphics[width=0.75\textwidth]{../lectures/figures/segmentation.png}
\caption{Foreground/background segmentation. \emph{From
\href{https://ai.stanford.edu/~syyeung/cvweb/tutorials.html}{https://ai.stanford.edu/\textasciitilde syyeung/cvweb/tutorial3.html}}}
\end{figure}


\end{frame}

\begin{frame}{Motivation}\label{motivation}
\begin{figure}
\centering
\includegraphics[width=0.6\textwidth]{../lectures/figures/kde.png}
\caption{Kernel density estimate. \emph{From
\href{https://scikit-learn.org/stable/auto_examples/neighbors/plot_species_kde.html}{Scikit-learn
KDE Demo}.}}
\end{figure}
\end{frame}


\begin{frame}[t, allowframebreaks]{Mixture Models}\label{mixture-models}

Let, 
\begin{itemize}
    \item \(N\) denote the number of data points 
    \item \(K\) denote the number of mixture components (i.e., clusters) 
    \item \(\mbx_n \in \reals^D\) denote the \(n\)-th data point 
    \item \(z_n \in \{1, \ldots, K\}\) be a latent variable denoting the cluster assignment of the \(n\)-th data point
\end{itemize}

The model is parameterized by, 
\begin{itemize}
    \item \(\mbtheta_k\), the natural parameters of
cluster \(k\). 
    \item \(\mbpi \in \Delta_{K-1}\), the cluster proportions (probabilities).
\end{itemize}
\end{frame}

\begin{frame}[t, allowframebreaks]{Mixture Models}\label{mixture-models}
The generative model is as follows:

\begin{enumerate}
\item
  Sample the assignment of each data point: \begin{align*}
  z_n &\iid{\sim} \mathrm{Cat}(\mbpi) \quad \text{for } n = 1, \ldots, N
  \end{align*}
\item
  Sample data points given their assignments: \begin{align*}
  \mbx_n &\sim p(\mbx \mid \mbtheta_{z_n}) \quad \text{for } n = 1, \ldots, N
  \end{align*}
\end{enumerate}
\end{frame}

\begin{frame}[t, allowframebreaks]{Joint distribution}\label{joint-distribution}

The joint distribution of the data and latent variables is,

\begin{align*}
    p(\{(z_n, \mbx_n)\}_{n=1}^N) 
    &= \prod_{n=1}^N \mathrm{Cat}(z_n \mid \pi) \, p(\mbx_n \mid z_n) \\
    &= \prod_{n=1}^N \prod_{k=1}^K \left[ \pi_k \, p(\mbx_n \mid \mbtheta_k) \right]^{\bbI[z_n = k]}
\end{align*}
\end{frame}

\begin{frame}[t, allowframebreaks]{Exponential family mixture
models}\label{exponential-family-mixture-models}

Assume an exponential family likelihood of the form, \begin{align*}
    p(\mbx \mid \mbeta_k) &= h(\mbx_n) \exp \left \{\langle t(\mbx_n), \mbeta_k \rangle - A(\mbeta_k) \right \}
\end{align*}

\textbf{Example:} Gaussian Mixture Model (GMM)

Assume the conditional distribution of \(\mbx_n\) is a Gaussian with
mean \(\mbtheta_k \in \reals^D\) and identity covariance: \begin{align*}
    p(\mbx_n \mid \mbtheta_k) &= \mathrm{N}(\mbx_n \mid \mbtheta_k, \mbI)
\end{align*}

Since we are assuming identity covariance, the sufficient statistics of the Gaussian
are \(t(\mbx) = \mbx\) and the natural parameters are simply \(\mbeta_k = \mbtheta_k\).
\end{frame}


\begin{frame}[t, allowframebreaks]{Two Inference Algorithms}\label{two-inference-algorithms}

Let's stick with the Gaussian mixture model for now. Suppose we observe
data points \(\{\mbx_n\}_{n=1}^N\) and want to infer the assignments
\(\{z_n\}_{n=1}^N\) and estimate the means \(\{\mbtheta_k\}_{k=1}^K\).
Here are two intuitive algorithms.
\end{frame}

\begin{frame}[t, allowframebreaks]{K-Means}\label{map-inference-and-k-means}

Suppose we knew the cluster assignments, \(z_n\). Then it would be
straightforward to estimate the cluster means: we could simply use the
maximum likelihood estimate,
\(\hat{\mbtheta}_{\mathsf{MLE}} = \frac{1}{N_k} \sum_{n: z_n=k} \mbx_n\),
where \(N_k = \sum_n \mathbb{I}[z_n =k]\) is the number of data poitns
in cluster \(k\).

Likewise, if we knew the cluster means, it would be straightforward to
compute the \emph{maximum a posteriori} (MAP) estimate of \(z_n\): we
would assign each data point to the nearest cluster. If we alternate
these two steps, we obtain the \textbf{K-Means algorithm}:
\end{frame}

\begin{frame}[t, allowframebreaks]{K-Means}\label{map-inference-and-k-means}
\textbf{Repeat} until convergence: 

\begin{enumerate}
    \item For each \(n=1,\ldots, N\), fix the means \(\mbtheta\) and set,
    \begin{align*}
        z_n &= \hat{z}_{n,\mathsf{MAP}} \\
        &= \arg \max_{k \in \{1,\ldots,K\}} \mathrm{N}(\mbx_n \mid \mbtheta_k, \mbI) \\
        &= \arg \min_{k \in \{1,\ldots, K\}} \|\mbx_n - \mbtheta_k\|_2
    \end{align*}
    
    \item For each \(k=1,\ldots,K\), fix all assignments \(\mbz\) and set,
    \begin{align*}
       \mbtheta_k &= \hat{\mbtheta}_{\mathsf{MLE}} \\
       &= \frac{1}{N_k} \sum_{n=1}^K \bbI[z_n=k] \mbx_n
    \end{align*}
\end{enumerate}
\end{frame}

\begin{frame}[t]{Comments on K-Means}

\begin{itemize}    
    \item \textbf{Question:} What does this algorithm implicitly assume
about the cluster probabilities \(\mbpi\)? How would you modify this
algorithm to incorporate and estimate \(\mbpi\)?

\item \textbf{Connection between K-Means and MAP estimation} 
Note that if we put an improper uniform prior on \(\mbtheta_k\), we could
think of this entire algorithm as MAP estimation of \(\mbtheta\) and
\(\mbz\) via \textbf{coordinate ascent}! 

\item K-Means made \textbf{hard assignments} of data points to clusters in
each iteration. That sounds a little extreme --- do you really want to
attribute a datapoint to a single class when it is right in the middle
of two clusters? What could we do instead?
\end{itemize}

\end{frame}

\begin{frame}[t, allowframebreaks]{Expectation Maximization (EM) for a GMM}\label{maximum-likelihood-estimation-via-em}
\textbf{Repeat} until convergence: 
\begin{enumerate}
    \item For each data point \(n\) and component \(k\), compute the
    \textbf{responsibility}: 
    \begin{align*}
        \omega_{nk} = \frac{\pi_k \mathrm{N}(\mbx_n \mid \mbtheta_k, \mbI)}{\sum_{j=1}^K \pi_j \mathrm{N}(\mbx_n \mid \mbtheta_j, \mbI)}
    \end{align*}

    \item
    For each component \(k\), update the mean: \begin{align*}
       \mbtheta_k^\star &= \frac{1}{N_k} \sum_{n=1}^K \omega_{nk} \mbx_n
    \end{align*}
\end{enumerate}

This is the \textbf{Expectation-Maximization (EM) algorithm}. As we will
show, EM yields an estimate that maximizes the \emph{marginal}
likelihood of the data.
\end{frame}

\begin{frame}[t, allowframebreaks]{Theoretical Motivation for EM}\label{theoretical-motivation}

Rather than maximizing the \textbf{joint probability}, EM is maximizing
the \textbf{marginal probability}, \begin{align*}
    \log p(\{\mbx_n\}_{n=1}^N; \mbtheta) 
    &= \log \sum_{z_1=1}^K \cdots \sum_{z_N=1}^K p(\{\mbx_n, z_n\}_{n=1}^N ; \mbtheta) \\
    &= \log \prod_{n=1}^N \sum_{z_n=1}^K p(\mbx_n, z_n ; \mbtheta) \\
    &= \sum_{n=1}^N  \log \sum_{z_n=1}^K p(\mbx_n, z_n ; \mbtheta) 
\end{align*} For discrete mixtures (with small enough \(K\)) we can
evaluate the log marginal probability. We can usually evaluate its
gradient too, so we could just do gradient ascent to find
\(\mbtheta^*\). However, EM typically obtains faster convergence rates.
\end{frame}

\begin{frame}[t, allowframebreaks]{The ELBO from another angle}\label{evidence-lower-bound-elbo}

The key idea is to obtain a lower bound on the marginal probability,
\begin{align*}
    \log p(\{\mbx_n\}_{n=1}^N; \mbtheta) 
    &= \sum_{n=1}^N  \log \sum_{z_n} p(\mbx_n, z_n; \mbtheta) \\
    &= \sum_{n=1}^N  \log \sum_{z_n} q(z_n) \frac{p(\mbx_n, z_n; \mbtheta)}{q(z_n)} \\
    &= \sum_{n=1}^N  \log \E_{q(z_n)} \left[\frac{p(\mbx_n, z_n; \mbtheta)}{q(z_n)} \right]
\end{align*} where \(q(z_n)\) is any distribution on
\(z_n \in \{1,\ldots,K\}\) such that \(q(z_n)\) is \textbf{absolutely
continuous} w.r.t. \(p(\mbx_n, z_n; \mbtheta)\).

\textbf{Jensen's Inequality:} states that, 
\begin{align*}
    f(\E[Y]) \geq \E\left[ f(Y) \right]
\end{align*} if \(f\) is a \textbf{concave function}, with equality iff
\(f\) is linear.

Applied to the log marginal probability, Jensen's inequality yields,
\begin{align*}
    \log p(\{\mbx_n\}_{n=1}^N; \mbtheta) 
    &= \sum_{n=1}^N  \log \E_{q_n} \left[\frac{p(\mbx_n, z_n; \mbtheta)}{q_n(z_n)} \right] \\
    &\geq \sum_{n=1}^N  \E_{q_n} \left[\log p(\mbx_n, z_n; \mbtheta) - \log q_n(z_n) \right] \\
    &\triangleq \cL[\mbtheta, \mbq]
\end{align*} where \(\mbq = (q_1, \ldots, q_N)\) is a tuple of
distributions, one for each latent variable \(z_n\).

This is called the \textbf{evidence lower bound}, or \textbf{ELBO} for
short. It is a function of \(\mbtheta\) and a \emph{functional} of
\(\mbq\), since each \(q_n\) is a probability density function. We can
think of \emph{EM as coordinate ascent on the ELBO}, alternating between
updating the parameters \(\mbtheta\) and the posteriors \$q.
\end{frame}

\begin{frame}[t, allowframebreaks]{M-step: Gaussian case}\label{m-step-gaussian-case}

Suppose we fix \(\mbq\). Since each \(z_n\) is a discrete latent
variable, \(q_n\) must be a probability mass function. Let it be denoted
by, \begin{align*}
    q_n &= [\omega_{n1}, \ldots, \omega_{nK}]^\top.
\end{align*} (These will be the \textbf{responsibilities} from before.)

Note that $q_n(k) = \Pr(z_n = k) = \mathbb{E}[\bbI[z_n=k]] = \omega_{nk}$ is the probability that the $n$-th data point belongs to cluster $k$.

Now, recall our basic model,
\(\mbx_n \sim \mathrm{N}(\mbtheta_{z_n}, \mbI)\), \begin{align*}
    \cL[\mbtheta, \mbq] 
    &= \sum_{n=1}^N \E_{q_n} [\log p(\mbx_n, z_n ; \mbtheta)] + c \\
    &= \sum_{n=1}^N \E_{q_n} [\bbI[z_n = k] \log p(\mbx_n, \mbtheta_k)] + c \\
    &= \sum_{n=1}^N \sum_{k=1}^K \omega_{nk} \log p(\mbx_n, \mbtheta_k) + c \\
    &= \sum_{n=1}^N \sum_{k=1}^K \omega_{nk} \left[ \mbx_n^\top \mbtheta_k - \tfrac{1}{2} \mbtheta_k^\top \mbtheta_k \right] + c
\end{align*}

Zooming in on just \(\mbtheta_k\), \begin{align*}
    \cL[\mbtheta, \mbq] 
    &= {\mbh}_{k}^\top \mbtheta_k - \tfrac{1}{2} \mbtheta_k^\top \mbJ_k \mbtheta_k
\end{align*} where \begin{align*}
    {\mbh}_{k} &= \sum_{n=1}^N \omega_{nk} \mbx_n
    \qquad
    {\mbJ}_{k} = \sum_{n=1}^N \omega_{nk} \mbI
\end{align*} Taking derivatives and setting to zero yields,
\begin{align*}
    \mbtheta_k^\star &=  \mbJ_k^{-1} \mbh_k 
    = \frac{\sum_{n=1}^N \omega_{nk} \mbx_n}{\sum_{n=1}^N \omega_{nk}}.
\end{align*} These are the same as the EM updates shown above!
\end{frame}

\begin{frame}[t, allowframebreaks]{E-step: Gaussian case}\label{e-step-gaussian-case}

As a function of \(q_n\), for discrete Gaussian mixtures with identity
covariance, \begin{align*}
    \cL[\mbtheta, \mbq] 
    &= \E_{q_n}\left[ \log p(\mbx_n, z_n; \mbtheta) - \log q_n(z_n)\right] + c \\
    &= \sum_{k=1}^K \omega_{nk} \left[ \log \mathrm{N}(\mbx_n \mid \mbtheta_k, \mbI) + \log \pi_k - \log \omega_{nk} \right] + c
\end{align*} where \(\mbpi = [\pi_1, \ldots, \pi_K]^\top\) is the vector
of prior cluster probabilities.

We also have two constraints: \(\omega_{nk} \geq 0\) and
\(\sum_k \omega_{nk} = 1\). Let's ignore the non-negative constraint for
now (it will automatically be satisfied anyway) and write the Lagrangian
with the simplex constraint, \begin{align*}
    \cJ(\mbomega_n, \lambda) &= \sum_{k=1}^K \omega_{nk} \left[ \log \mathrm{N}(\mbx_n \mid \theta_k, \mbI) + \log \pi_k - \log \omega_{nk} \right] - \lambda \left( 1 - \sum_{k=1}^K \omega_{nk} \right)
\end{align*}

Taking the partial derivative wrt \(\omega_{nk}\) and setting to zero
yields, \begin{align*}
    \frac{\partial}{\partial \omega_{nk}} \cJ(\mbomega_n, \lambda) 
    &= \log \mathrm{N}(\mbx_n \mid \mbtheta_k, \mbI) + \log \pi_k - \log \omega_{nk} - 1 + \lambda = 0 \\
    \Rightarrow \log \omega_{nk}^\star &= \log \mathrm{N}(\mbx_n \mid \mbtheta_k, \mbI) + \log \pi_k + \lambda - 1 \\
    \Rightarrow \omega_{nk}^\star &\propto \pi_k \mathrm{N}(\mbx_n \mid \mbtheta_k, \mbI) 
\end{align*} Enforcing the simplex constraint yields, \begin{align*}
    \omega_{nk}^\star &= \frac{\pi_k \mathrm{N}(\mbx_n \mid \mbtheta_k, \mbI)}{\sum_{j=1}^K \pi_j \mathrm{N}(\mbx_n \mid \mbtheta_j, \mbI)},
\end{align*} just like above.

Note that \begin{align*}
    \omega_{nk}^\star &\propto p(z_n=k) \, p(\mbx_n \mid z_n=k, \mbtheta) = p(z_n = k \mid \mbx_n, \mbtheta) .
\end{align*} That is, the responsibilities equal the posterior
probabilities!
\end{frame}

\begin{frame}[t, allowframebreaks]{The ELBO is tight after the
E-step}\label{the-elbo-is-tight-after-the-e-step}

Equivalently, \(q_n\) equals the posterior,
\(p(z_n \mid \mbx_n, \mbtheta)\). At that point, the ELBO simplifies to,
\begin{align*}
    \cL[\mbtheta, \mbq] 
    &= \sum_{n=1}^N  \E_{q_n} \left[\log p(\mbx_n, z_n ; \mbtheta) - \log q_n(z_n) \right] \\
    &= \sum_{n=1}^N  \E_{p(z_n \mid \mbx_n, \mbtheta)} \left[\log p(\mbx_n, z_n ; \mbtheta) - \log p(z_n \mid \mbx_n; \mbtheta) \right] \\
    &= \sum_{n=1}^N  \E_{p(z_n \mid \mbx_n; \mbtheta)} \left[\log p(\mbx_n ; \mbtheta) \right] \\
    &= \sum_{n=1}^N \log p(\mbx_n ; \mbtheta) \\
    &= \log p(\{\mbx_n\}_{n=1}^N, \mbtheta)
\end{align*}


We can view the EM algorihtm as a \textbf{minorize-maximize (MM)}
algorithm where we iteratively lower bound the ELBO and and then
maximize the lower bound.
\end{frame}


\begin{frame}[t, allowframebreaks]{M-step: General Case}\label{m-step-general-case}

Now let's consider the general Bayesian mixture with exponential family
likelihoods.

While we're at it, let's also add conjugate priors \(p(\mbtheta)\) with
pseudo-counts \(\nu\) and pseudo-observations \(\mbchi\). As a function
of \(\mbtheta\), \begin{align*}
    \cL[\mbtheta, \mbq] 
    &= \log p(\mbtheta) + 
    \sum_{n=1}^N \E_{q_n} [\log p(\mbx_n, z_n \mid \mbtheta)] + \mathrm{c} \\
    &= \log p(\mbtheta) + 
    \sum_{n=1}^N \sum_{k=1}^K \omega_{nk} \log p(\mbx_n \mid \mbtheta_k) + \mathrm{c} \\
    &= \sum_{k=1}^K \left[\mbchi^\top \mbtheta_k - \nu A(\mbtheta_k)\right] +
    \sum_{n=1}^N \sum_{k=1}^K \omega_{nk} \left[ t(\mbx_n)^\top \mbtheta_k - A(\mbtheta_k) \right] + \mathrm{c} 
\end{align*}

Zooming in on just \(\mbtheta_k\), \begin{align*}
    \cL[\mbtheta, \mbq] 
    &= \mbchi_k'^\top \mbtheta_k - \nu_k' A(\mbtheta_k)
\end{align*} where \begin{align*}
    \mbchi_k' &= \mbchi + \sum_{n=1}^N \omega_{nk} t(\mbx_n)
    \qquad
    \nu_k' = \nu + \sum_{n=1}^N \omega_{nk}
\end{align*} Taking derivatives and setting to zero yields,
\begin{align*}
    \label{eq:gen_mstep}
    \mbtheta_k^* &= \left[\nabla A \right]^{-1} \left(\frac{\mbchi_k'}{\nu_k'}\right)
\end{align*}

Recall that \(\nabla A^{-1}: \cM \mapsto \Omega\) is a mapping from mean
parameters to natural parameters (and the inverse exists for minimal
exponential families). Thus, the generic M-step above amounts to finding
the natural parameters \(\mbtheta_k^*\) that yield the expected
sufficient statistics \(\mbchi_{k}' / \nu_k'\) by inverting the gradient
mapping.
\end{frame}

\begin{frame}[t, allowframebreaks]{E-step: General Case}\label{e-step-general-case}

In our first pass, we assumed \(q_n\) was a finite pmf. More generally,
\(q_n\) will be a probability density function, and optimizing over
functions usually requires the \emph{calculus of variations}. (Ugh!)

However, note that we can write the ELBO in a slightly different form,
\begin{align*}
    \cL[\mbtheta, \mbq] 
    &= \log p(\mbtheta) + \sum_{n=1}^N  \E_{q_n} \left[\log p(\mbx_n, z_n \mid \mbtheta) - \log q_n(z_n) \right] \\
    &= \log p(\mbtheta) + \sum_{n=1}^N  \E_{q_n} \left[\log p(z_n \mid \mbx_n, \mbtheta) + \log p(\mbx_n \mid \mbtheta) - \log q_n(z_n) \right] \\
    &= \log p(\mbtheta) + \sum_{n=1}^N  \left[\log p(\mbx_n \mid \mbtheta) -\KL{q_n(z_n)}{p(z_n \mid \mbx_n, \mbtheta)} \right] \\
    &= \log p(\{\mbx_n\}_{n=1}^N, \mbtheta) - \sum_{n=1}^N  \KL{q_n(z_n)}{p(z_n \mid \mbx_n, \mbtheta)}
\end{align*} where \(\KL{\cdot}{\cdot}\) denote the
\textbf{Kullback-Leibler divergence}. (Note that we included the prior
on \(\mbtheta\) since we are treating it as a random variable with a
prior in this general case.)

Recall, the KL divergence is defined as, \begin{align*}
    \KL{q(z)}{p(z)} &= \int q(z) \log \frac{q(z)}{p(z)} \dif z.
\end{align*}

It gives a notion of how similar two distributions are, but it is
\emph{not a metric!} (It is not symmetric.) Still, it has some intuitive
properties: 1. It is non-negative, \(\KL{q(z)}{p(z)} \geq 0\). 2. It
equals zero iff the distributions are the same,
\(\KL{q(z)}{p(z)} = 0 \iff q(z) = p(z)\) almost everywhere.

Maximizing the ELBO wrt \(q_n\) amounts to minimizing the KL divergence
to the posterior \(p(z_n \mid \mbx_n, \mbtheta)\), \begin{align*}
    \cL[\mbtheta, \mbq] 
    &= \log p(\mbtheta) + \sum_{n=1}^N  \left[\log p(\mbx_n \mid \mbtheta) -\KL{q_n(z_n)}{p(z_n \mid \mbx_n, \mbtheta)} \right] \\
    &= -\KL{q_n(z_n)}{p(z_n \mid \mbx_n, \mbtheta)} + c
\end{align*} As we said, the KL is minimized when
\(q_n(z_n) = p(z_n \mid \mbx_n, \mbtheta)\), so the optimal update is,
\begin{align*}
    q_n^\star(z_n) &= p(z_n \mid \mbx_n, \mbtheta),
\end{align*} just like we found above.
\end{frame}

\begin{frame}[t, allowframebreaks]{Conclusion}\label{conclusion}

Mixture models are basic building blocks of statistics, and our first
encounter with \textbf{discrete latent variable models (LVMs)}. (Where
have we seen continuous LVMs so far?) Mixture models have widespread
uses in both density estimation (e.g., kernel density estimators) and
data science (e.g., clustering). Next, we'll talk about how to extend
mixture models to cases where the cluster assignments are correlated in
time.
\end{frame}

    
    
\end{document}
